\chapter{系统分析与设计}

\section{系统概述}

本文设计的企业供应链协同管理系统面向制造企业订单履约、库存控制、生产执行、原材料采购、供应商协同与质量追溯等关键业务场景，目标是在统一平台内建立覆盖顾客、销售管理员、仓库管理员、生产管理员、采购管理员、供应商、质检管理员以及系统管理员的协同工作机制。与仅面向单一部门的传统业务系统相比，本文系统更强调跨角色、跨状态和跨环节的数据联动，即以订单对象为主线，以库存为中介，以生产与采购为补充路径，以质检与追溯为质量保障手段，形成较为完整的业务闭环。

从系统边界看，本文实现的系统既包括企业内部管理流程，也包含供应商侧的外部协同功能。顾客能够提交订单并查看本人订单处理结果；销售管理员负责订单入口审核与销售记录维护；仓库管理员承担库存核查、出入库、发货及库存预警处理；生产管理员依据订单缺口或计划建议执行生产；采购管理员面向原材料短缺场景创建采购订单并跟踪供应商反馈；供应商根据自身账号查看与自己相关的原材料档案和采购订单；质检管理员对生产批次进行检验并根据结果决定是否退回生产或进入仓储入库流程；系统管理员则负责内部账号维护与系统级监督。由此可见，系统的核心不在于单一业务模块的实现，而在于通过统一的数据模型和统一的状态管理机制支撑多角色协同。

结合当前项目实现情况，系统采用前后端分离架构。前端基于 Vue 3 构建角色化业务页面，后端基于 Spring Boot 提供统一 REST 接口，并通过 Spring Security 与 JWT 完成身份认证和访问控制，通过 WebSocket/STOMP 完成关键业务事件的实时推送。数据库层使用 MySQL 存储结构化业务数据，JPA 负责实体映射与持久化访问。该总体方案使系统既能满足功能实现需求，也具备较好的工程可维护性与扩展性。

\section{需求分析}

\subsection{用户角色分析}

供应链协同系统首先面对的是多角色参与带来的职责划分问题。不同角色虽然操作对象可能相同，例如均围绕订单、产品、库存或批次开展工作，但其可执行动作和责任范围存在明显差异。如果在需求分析阶段不能对角色职责进行清晰界定，则在系统设计时容易出现权限交叉、流程混乱和责任不明确的问题。因此，本文首先从实际业务出发对各类角色进行分析。

\begin{table}[htbp]
\centering
\footnotesize
\caption{系统主要用户角色及职责分析}
\label{tab:roles-analysis-real}
\begin{tabular}{|c|p{3.4cm}|p{7.8cm}|}
\hline
角色 & 业务目标 & 主要职责与权限边界 \\\hline
系统管理员 & 保障系统基础运行与组织管理 & 维护内部人员账号、分配角色、启停账号、查看全局业务记录；不直接参与订单审核、采购执行或库存作业，但对系统运行结果负有监督职责。 \\\hline
顾客 & 提交购买需求并跟踪结果 & 进行注册登录、下达订单、查看本人订单明细与金额信息、根据订单号查询质量追溯结果；无权进入企业内部管理模块。 \\\hline
销售管理员 & 控制订单入口质量 & 查看顾客下单请求，执行接受或拒绝处理，维护销售记录，推进订单进入仓库核查；不负责库存出入库、采购执行和账号管理。 \\\hline
仓库管理员 & 保障库存真实与履约执行 & 查看成品与原材料库存，执行入库、出库、发货、库存流水管理与库存预警处理；在订单流程中负责库存核查、预留库存、发货和生产入库确认。 \\\hline
生产管理员 & 保障成品供给能力 & 接收生产计划订单，执行生产任务与完工回传，查看自身生产记录和质检退回信息；不能直接越过质检环节入库。 \\\hline
采购管理员 & 保障原材料供应连续性 & 根据库存预警、周计划或人工需求创建采购单，跟踪供应商反馈，通知仓库收货；不承担库存管理职责。 \\\hline
供应商 & 响应企业采购协同 & 查看与自身绑定的原材料档案，处理采购单的接单、拒单和发货操作，查看自身供应记录；不能访问企业内部库存和订单管理功能。 \\\hline
质检管理员 & 控制成品质量与异常闭环 & 对生产完成后的批次进行合格/不合格判定，记录质检原因并在不合格时通知对应生产管理员；仅质检合格的批次才能进入仓库入库确认阶段。 \\\hline
\end{tabular}
\end{table}

由表~\ref{tab:roles-analysis-real} 可知，系统角色之间并非简单的页面访问差异，而是对应完整的业务职责边界。例如，销售管理员和仓库管理员都可以查看订单，但前者关注审核与销售记录，后者关注库存满足性与发货执行；生产管理员能够结束生产，但无法绕过质检直接变更库存；采购管理员能够创建采购订单，却不能直接执行库存入库。正因为角色职责边界清晰，系统才有必要在接口级、页面级和服务级同时实现权限控制。

\subsection{功能需求分析}

结合项目实际实现情况，本文系统的功能需求可归纳为以下几个方面。

\subsubsection{认证授权与账号管理需求}

系统应支持多角色用户的统一注册、登录与身份识别，并在登录后根据角色范围开放不同功能入口。系统管理员还需具备内部人员账号管理能力，包括查看账号列表、创建内部账号、修改角色以及启用或禁用账号等。由于系统采用前后端分离模式，该模块不仅要满足登录校验需求，还要满足后续接口调用中无状态身份识别需求。

\subsubsection{顾客订单与销售审核需求}

顾客侧应具备下单、查看个人订单和质量追溯查询功能。销售管理员需要查看待审核订单，并对订单执行接受或拒绝处理；当销售接受后，订单需自动进入仓库核查阶段；若销售拒绝，系统应保留拒绝理由并终止后续流转。由于系统同时存在多个销售管理员，订单审核必须保证一次有效处理，避免重复审核。

\subsubsection{仓储与库存管理需求}

仓库管理员需要查看成品和原材料库存明细，包括产品名称、库存数量、预留数量、仓库信息以及入库/出库流水记录。对订单而言，仓库管理员需在接到已审核订单后判断库存是否充足；库存充足时执行预留和发货，不充足时需生成生产计划。对采购到货和生产完工场景，仓库管理员还需承担最终入库确认职责。系统还应具备库存预警功能，在成品或原材料低于安全库存时向对应角色发送提示。

\subsubsection{生产协同需求}

当订单库存不足时，系统应能自动生成生产计划订单并通知生产管理员。生产管理员需要查看待执行计划、提交完工结果并形成生产记录。项目实现中，生产完工后并不直接自动入库，而是先生成待检批次并通知质检管理员，只有质检合格后仓库才能执行确认入库。因此，生产模块不仅要完成任务登记，还要与订单状态、批次生成和后续质检形成联动。

\subsubsection{采购与供应商协同需求}

采购管理员可以主动创建采购单，也可以根据库存预警或周计划建议开展采购。采购单创建后需自动生成采购单号，并校验所选供应商是否与原材料档案中维护的供应商信息一致。供应商收到采购通知后可执行接受、拒绝或发货，采购管理员在供应商发货后通知仓库收货，仓库收货确认后系统自动完成原材料入库并生成库存流水。该过程中的状态变化需要具备可追踪性。

\subsubsection{质检与质量追溯需求}

系统需支持以生产批次为核心的质检管理机制。生产完成后生成批次记录，质检管理员对该批次进行合格或不合格判定；若检验不合格，系统通知对应生产管理员返工并将订单退回生产阶段；若检验合格，则通知仓库执行后续入库确认。与此同时，顾客可基于订单号查询对应订单是否存在质检记录及其检验结果，以满足质量追溯需求。

\subsubsection{库存预警与周计划需求}

系统应具备面向成品与原材料的库存预警能力。对于成品，当可用库存不足时，仓库管理员可以根据预警结果创建生产计划；对于原材料，当库存低于安全库存时，可生成采购申请或采购建议。除此之外，系统还需按周生成生产周计划与采购周计划：生产周计划综合考虑上周完工量、当前未完成订单需求和当前可用成品库存；采购周计划综合考虑 BOM 需求、上周采购量、当前库存、在途采购量和安全库存缺口。

\subsubsection{实时消息通知需求}

由于系统涉及多个角色围绕同一业务对象协同处理，实时通知是提升处理效率的重要手段。系统应在订单流转、生产任务生成、质检完成、采购状态变化、库存预警触发和周计划生成等关键节点主动推送消息，从而减少人工催办和轮询刷新带来的时间损耗。

\subsection{非功能需求分析}

除功能需求外，系统还需满足若干非功能性要求。

\begin{enumerate}
  \item \textbf{安全性要求。} 系统必须具备身份认证、角色鉴权和关键操作校验能力，避免未授权用户访问敏感接口，避免通过伪造请求绕过前端限制。
  \item \textbf{一致性要求。} 订单状态、库存数量、生产计划、批次状态和采购订单状态需保持一致，关键业务更新应在统一事务中完成，避免出现部分成功、部分失败的中间状态。
  \item \textbf{可追溯性要求。} 销售记录、库存流水、生产记录、采购记录和质检记录都需要保留操作人信息、时间信息与业务关联信息，以支撑责任定位与质量追溯。
  \item \textbf{实时性要求。} 对于订单受理、仓库核查、生产完工、质检结果、供应商发货等时效性较强场景，系统应提供及时消息同步能力。
  \item \textbf{可维护性与扩展性要求。} 系统应采用清晰的分层结构与模块边界，以便后续扩展统计分析、移动端接入、更多计划算法或细粒度权限策略。
  \item \textbf{易用性要求。} 不同岗位用户应能基于自身职责快速定位操作页面，系统页面结构应尽量与业务流程一致，减少学习成本。
\end{enumerate}

\section{整体架构设计}

\subsection{系统总体架构}

从实现层面看，系统总体架构可划分为前端表示层、接口控制层、业务服务层、数据持久层和消息通信层五个部分。

\begin{table}[htbp]
\centering
\footnotesize
\caption{系统总体架构层次说明}
\label{tab:overall-arch-layers}
\begin{tabular}{|c|p{3.4cm}|p{7.8cm}|}
\hline
层次 & 主要技术/组件 & 主要职责 \\\hline
前端表示层 & Vue 3、Vite、Vue Router、Axios & 提供登录、订单、库存、采购、生产、质检、账号管理等页面；根据用户角色动态展示菜单与待办信息。 \\\hline
接口控制层 & AuthController、SalesOrderController、InventoryController、ProcurementController、ProductionController、QualityController、AdminController 等 & 接收前端请求、完成参数校验和角色访问入口控制，对外暴露统一的 REST 接口。 \\\hline
业务服务层 & OrderWorkflowService、ProcurementWorkflowService、QualityService、WeeklyPlanningService 等 & 封装订单流转、采购协同、质检回流、库存预警、周计划生成等核心业务规则。 \\\hline
数据持久层 & Spring Data JPA、Repository 接口、MySQL & 管理用户、订单、库存、批次、采购单、周计划等结构化数据，实现实体映射与持久化访问。 \\\hline
消息通信层 & WebSocket、SockJS、STOMP、NotificationService & 对订单、生产、采购、质检、库存预警等事件进行实时广播，支撑多角色消息同步。 \\\hline
\end{tabular}
\end{table}

其中，前端与后端之间主要通过 HTTP 接口进行数据交互，而关键业务事件则通过 WebSocket 通道主动推送。该架构使系统既具备面向事务处理的标准化接口，也具备面向协同提醒的实时通信能力。对于供应链协同系统而言，这种“接口调用 + 事件通知”的复合机制能够更好地适配订单驱动型业务场景。

\subsection{前后端分层设计}

后端分层遵循控制层、服务层与持久层分离原则。控制层不直接承担复杂业务计算，而是将核心处理委托给服务层，从而保证业务规则集中管理。例如，订单由销售审核进入仓库核查、仓库判断库存是否满足、库存不足时生成生产计划、生产完成后转入待质检状态、质检合格后转待生产入库等动作，均由服务层统一完成状态控制。这样既提高了逻辑可维护性，也保证了不同接口入口下的规则一致性。

前端部分则采用页面层、接口封装层与公共组件层相结合的方式。项目中通过统一的接口服务文件对订单、生产、库存、采购、质检、产品、客户和管理员功能进行封装，页面再基于这些接口完成视图展示和交互控制。由于同一系统需要服务多类角色，因此前端设计重点不在于单页面复杂度，而在于如何让不同角色以最少操作完成本岗位工作。

\subsection{权限控制设计}

权限控制设计围绕“角色最小授权”原则展开。系统后端通过 Spring Security 在请求入口执行认证与路径授权，在业务服务层再执行角色和状态双重校验。以当前项目实现为例，\texttt{/api/v1/admin/**} 仅允许系统管理员访问，\texttt{/api/v1/customer/**} 主要面向顾客与管理员，\texttt{/api/v1/orders/**} 面向销售、仓库、生产和管理员，\texttt{/api/v1/inventory/**} 主要面向仓库管理员与管理员，\texttt{/api/v1/procurement/**} 面向采购管理员、供应商、仓库管理员及管理员，\texttt{/api/v1/quality/**} 面向质检管理员与管理员。

在服务层中，权限控制并未因为页面隐藏而放松。例如，即使用户通过构造请求试图执行并非自身职责范围内的动作，后端仍会通过当前认证信息和目标对象状态进行二次校验。以订单流程为例，只有销售管理员才能将顾客订单由“待销售审核”推进到“待仓库核查”；只有仓库管理员才能执行发货；只有生产管理员才能回传生产完工；只有质检管理员才能对批次执行最终检验。这种“页面控制 + 接口控制 + 服务校验”的三层机制能够有效降低越权风险。

\section{业务流程设计}

\subsection{客户下单与销售审核流程}

顾客通过前端下单页面提交订单后，系统首先创建订单主记录和订单明细记录，并将订单状态初始化为待销售审核。销售管理员进入订单管理页面后，可查看订单编号、顾客信息、产品信息、数量、金额以及配送地址等内容，并根据业务规则执行接受或拒绝操作。

当销售管理员选择接受时，系统调用订单工作流服务将订单状态推进为待仓库核查，并向仓库管理员主题推送“接收一条新的订单 + 订单号 + 请查看”的通知；若销售管理员选择拒绝，则系统要求填写拒绝原因，并终止该订单的后续流程。由于项目中要求同一订单只能由一名销售管理员完成首次处理，因此该流程在服务层中还需校验订单当前状态是否仍处于可审核状态。

\subsection{仓储发货与库存流转流程}

仓库管理员接收已通过销售审核的订单后，首先对订单明细逐项检查库存是否满足发货需求。项目实现中，系统并非简单比较库存总量，而是通过 \texttt{quantity - reservedQuantity} 计算可用库存，并逐项判断是否存在产品缺口。若所有明细均满足，则系统预留对应库存，将订单状态更新为已接单，并通知销售侧订单已经完成库存核验；若存在缺口，则系统将订单状态更新为生产中，并根据缺口数量生成生产计划。

当订单处于已接单状态后，仓库管理员可以执行发货操作。系统在发货时优先消耗已预留库存，并同步扣减库存数量、更新预留数量、生成 \texttt{StockTransaction} 出库流水，并向顾客发送“您的订单 + 订单号 + 已发货”的通知。由此可见，仓库模块既承担库存检查职责，也承担订单履约执行职责，是连接销售、生产和顾客的重要枢纽。

\subsection{生产计划、生产执行与质检流程}

在库存不足场景下，系统会针对缺货产品生成生产计划记录 \texttt{ProductionPlan}。每条计划包含计划编号、产品、计划数量、开始时间、结束时间、创建人及状态等信息，并以订单号、产品编号和时间戳构成较为稳定的业务标识。生成计划后，系统向生产管理员主题广播生产通知。

生产管理员在完成生产后，并不会直接将成品写入库存，而是通过“生产完工回传”动作将计划状态更新为 \texttt{DONE}，同时为对应订单和产品生成待检批次 \texttt{Batch}。批次中保存批次号、来源订单号、产品、数量、生产管理员信息和当前质检状态等。随后，系统将订单状态推进为待质检，并通知质检管理员对该批次执行检验。

质检管理员对批次作出合格或不合格判定后，系统进入分支处理。若判定为不合格，则写入 \texttt{QualityRecord}，向对应生产管理员发送返工通知，并将订单状态退回生产中；若判定为合格，则更新批次状态为合格，将订单状态推进为待生产入库，并通知仓库管理员执行生产入库确认。最终由仓库管理员根据合格批次完成入库，系统再为订单预留可用库存并进入发货准备阶段。该流程体现了项目中“生产完成后通知仓库管理员，仓库在质检通过后确认入库”的真实业务要求。

\subsection{原材料采购与供应商协同流程}

原材料采购既可能由采购管理员主动发起，也可能由库存预警或采购周计划建议触发。采购管理员创建采购单时，系统自动生成采购单号 \texttt{poNo}，并校验所选供应商是否与原材料档案中的优选供应商一致，避免采购对象与档案信息不匹配。采购单创建成功后，状态进入“待供应商确认”，并向对应供应商主题发送采购通知。

供应商登录系统后，只能查看与自身账号相关的采购单和原材料档案。对于采购单，供应商可执行接单、拒单或发货操作。若接单，则采购状态变更为“供应商已接单”；若拒单，则采购状态终止并通知采购管理员；若发货，则采购状态变更为“供应商已发货”，采购管理员收到通知后可进一步通知仓库收货。仓库确认收货后，系统自动为采购明细对应原材料执行入库并生成库存流水，采购单状态变更为“已入库”。

\subsection{库存预警与周计划生成流程}

库存预警流程将成品与原材料分开处理。对于成品，系统基于产品安全库存、当前可用库存和在制量识别潜在缺口，并允许仓库管理员根据推荐数量创建生产计划；对于原材料，系统结合当前库存、在途采购量和安全库存缺口判断是否需要生成采购申请。这样一来，系统不仅能在订单发生时被动响应，也能在库存水平下降时主动发出业务信号。

在周计划生成方面，系统每周自动或手动计算生产周计划与采购周计划。生产周计划综合上周完工量、当前未完成订单需求和可用成品库存；采购周计划在此基础上继续结合 BOM 需求、上周采购量、当前可用原材料库存和在途采购量，形成对未来一周的采购建议。该流程将库存预警从单次事件扩展为面向周维度的计划支持机制。

\subsection{质量追溯流程}

质量追溯流程以订单号为入口，以批次和质检记录为核心。项目实现中，生产完工时生成的 \texttt{Batch} 保存了来源订单号；质检过程生成的 \texttt{QualityRecord} 又与批次和产品建立关联；仓库入库和发货过程产生的库存流水则记录了 lot、相关类型和相关业务编号。因此，当顾客在质量追溯页面输入订单号时，系统能够沿着“订单 $\rightarrow$ 批次 $\rightarrow$ 质检记录 $\rightarrow$ 入库/发货记录”的链路返回结果。

这一设计说明，质量追溯并不是独立模块临时拼接的结果，而是依赖业务全流程中的关键节点留痕。只有在设计阶段即明确数据关联方式，后续的追溯查询才具备可实现性与可信度。

\section{数据库设计}

\subsection{概念结构设计}

依据系统功能与业务流程，本文将数据库中的核心对象划分为用户与权限域、订单履约域、库存域、生产与质量域、采购域以及计划域六类。

\begin{itemize}
  \item \textbf{用户与权限域：} 包括 \texttt{User}、\texttt{Role} 以及中间关系表 \texttt{user\_roles}，用于描述系统账号、角色和授权关系。
  \item \textbf{订单履约域：} 包括 \texttt{SalesOrder}、\texttt{OrderItem}、\texttt{SalesRecord} 等实体，用于描述顾客订单、订单明细和销售处理结果。
  \item \textbf{库存域：} 包括 \texttt{Product}、\texttt{InventoryItem}、\texttt{StockTransaction}、\texttt{Warehouse} 等实体，用于描述库存基础信息、仓库维度数量和出入库流水。
  \item \textbf{生产与质量域：} 包括 \texttt{ProductionPlan}、\texttt{ProductionTask}、\texttt{Batch}、\texttt{QualityRecord} 等实体，用于描述生产计划、生产批次和质检结果。
  \item \textbf{采购域：} 包括 \texttt{PurchaseRequest}、\texttt{PurchaseOrder}、\texttt{PurchaseOrderItem} 等实体，用于描述原材料采购业务。
  \item \textbf{计划域：} 包括 \texttt{ProductionWeeklyPlan}、\texttt{ProductionWeeklyPlanItem}、\texttt{ProcurementWeeklyPlan}、\texttt{ProcurementWeeklyPlanItem}、\texttt{Bom} 和 \texttt{BomItem}，用于支撑周计划和物料需求换算。
\end{itemize}

这种概念划分方式既体现了实体之间的业务从属关系，也为后续服务层设计提供了稳定的数据边界。

\subsection{主要数据表设计}

与初版论文中采用的概括性表名不同，本文根据项目中的实际实体映射关系，对主要数据表进行更加贴近实现的说明，如表~\ref{tab:real-main-tables} 所示。

\begin{table}[htbp]
\centering
\footnotesize
\caption{系统主要数据表及关键字段设计}
\label{tab:real-main-tables}
\begin{tabular}{|c|p{4.0cm}|p{6.8cm}|}
\hline
表名 & 关键字段 & 作用说明 \\\hline
users & id、username、password、fullName、email、enabled、createdAt & 存储系统用户账号信息，是登录认证和角色绑定的基础。 \\\hline
roles & id、name & 存储角色定义，如 ROLE\_ADMIN、ROLE\_WAREHOUSE\_MANAGER 等。 \\\hline
user\_roles & user\_id、role\_id & 维护用户与角色之间的多对多关系。 \\\hline
products & id、sku、name、productType、preferredSupplier、unitPrice、safetyStock & 存储产品与原材料档案，其中 productType 区分成品与原材料。 \\\hline
sales\_order & id、orderNo、customer\_id、orderDate、status、totalAmount、shippingAddress & 存储顾客订单主信息及其当前处理状态。 \\\hline
order\_item & id、order\_id、product\_id、quantity、unitPrice、lineTotal & 存储订单明细项，实现一个订单对应多种产品。 \\\hline
inventory\_items & id、product\_id、warehouse\_id、quantity、reservedQuantity、lot、updatedAt & 存储仓库维度库存信息，reservedQuantity 用于订单预留。 \\\hline
stock\_transaction & id、transactionNo、product\_id、warehouse\_id、changeQuantity、transactionType、relatedType、relatedId、createdByName & 存储入库、出库、发货等库存流水，支撑审计和追溯。 \\\hline
production\_plan & id、planNo、product\_id、plannedQuantity、status、createdByName、completedByName & 存储订单缺口触发或预警触发的生产计划。 \\\hline
batch & id、batchNo、product\_id、quantity、sourceOrderNo、qualityStatus、productionManagerName & 存储生产完成后进入质检流程的批次信息，是质量追溯核心实体。 \\\hline
quality\_record & id、batch\_id、product\_id、inspectorName、inspectionDate、result、remarks & 存储质检记录与异常说明，并记录是否已通知生产管理员。 \\\hline
purchase\_order & id、poNo、supplier\_id、status、totalAmount、createdByName、shippedAt、receivedAt & 存储采购订单头信息及其供应链状态。 \\\hline
purchase\_order\_item & id、purchase\_order\_id、product\_id、quantity、unitPrice、lineTotal & 存储采购明细，描述具体原材料及数量。 \\\hline
production\_weekly\_plan & id、weekStart、weekEnd、generatedBy、algorithmNote、generatedAt & 存储生产周计划主记录。 \\\hline
procurement\_weekly\_plan & id、weekStart、weekEnd、generatedBy、algorithmNote、generatedAt & 存储采购周计划主记录。 \\\hline
bom / bom\_item & product\_id、componentProduct、quantity & 描述产品与原材料之间的物料清单关系，为采购周计划计算提供依据。 \\\hline
\end{tabular}
\end{table}

上述表结构体现了本系统数据库设计的两个特点。其一，系统在核心事务表之外增加了记录型表，例如库存流水表和质检记录表，以满足留痕与追溯需求；其二，系统引入周计划表和 BOM 表，说明其不仅是事务管理系统，也承担一定的业务决策支持职责。

\subsection{关键实体关系分析}

从实体关系角度看，系统的业务闭环主要通过以下几类关系加以体现。

\begin{enumerate}
  \item \textbf{用户与角色关系。} \texttt{User} 与 \texttt{Role} 为多对多关系，该设计可以兼容系统管理员、供应商、采购管理员等多角色用户的统一管理需求。
  \item \textbf{订单与明细关系。} \texttt{SalesOrder} 与 \texttt{OrderItem} 为一对多关系，一个订单可以包含多个产品明细，是订单金额计算和库存核查的基础。
  \item \textbf{产品与库存关系。} \texttt{Product} 与 \texttt{InventoryItem} 之间为一对多关系，原因在于同一产品可能分布于多个仓库；\texttt{InventoryItem} 再与 \texttt{Warehouse} 形成多对一关系。
  \item \textbf{订单与生产关系。} 当订单产生库存缺口时，\texttt{SalesOrder} 通过订单号间接关联到若干 \texttt{ProductionPlan}，而生产完工后又进一步生成 \texttt{Batch} 实体，从而将订单和质检链路连接起来。
  \item \textbf{批次与质检关系。} \texttt{Batch} 与 \texttt{QualityRecord} 为一对多关系，同一批次在返工后可能形成多次质检记录。该关系保证了质检过程具有历史可追溯性。
  \item \textbf{采购单与供应商关系。} \texttt{PurchaseOrder} 与 \texttt{User(ROLE\_SUPPLIER)} 为多对一关系，一个供应商可对应多张采购订单，而每张采购订单在当前实现中只对应一个供应商。
  \item \textbf{计划与明细关系。} 生产周计划与采购周计划均采用主表 + 明细表结构，用于存储每周计划总体信息和逐产品、逐原材料的建议数量信息。
\end{enumerate}

综合来看，实体关系设计并未停留在传统“订单、库存、采购”三个孤立对象上，而是围绕“订单驱动生产与采购、批次支撑质检与追溯、流水支撑库存审计、计划支撑决策建议”这一主线构建。这样的设计能够较好地契合项目的业务目标，并为后续系统详细实现和核心功能分析提供坚实的数据层基础。

